\subsection{ASYMPTOTIC EXPANSION OF THE PARTIAL PRODUCTS OF AN INFINITE PRODUCT}

One knows (Gen.~Top., V, p.~22 and 23) that for the infinite product with general factor \(1 + u_{n}\) ( \(u_{n} > -1\) ) to be convergent (resp. commutatively convergent) it is necessary and sufficient that the series with general term \(\log(1 + u_{n})\) should be convergent (resp. commutatively convergent), and that one then has the relation

\[
\log \prod_ {n = 1} ^ {\infty} (1 + u _ {n}) = \sum_ {n = 1} ^ {\infty} \log (1 + u _ {n}).
\]

When the infinite product converges one knows that \(u_{n}\) tends to 0; thus \(\log(1 + u_{n}) \sim u_{n}\) ; now one knows that for a series of real numbers to be commutatively convergent it is necessary and sufficient that it should be absolutely convergent (Gen.~Top., IV, p.~372, prop. 5); by prop. 1 one thus recovers the fact that the product with general factor \(1 + u_{n}\) is commutatively convergent if and only if the series with general term \(u_{n}\) is absolutely convergent (Gen.~Top., IV, p.~368, th. 4).

A similar argument applies to an infinite product whose general factor is a complex number \(1 + u_{n}\) ( \(u_{n} \neq -1\) ). Indeed, for such a product to be commutatively convergent it is necessary and sufficient (Gen.~Top., VIII, p.~115, prop. 2) that the infinite product with general factor \(|1 + u_{n}|\) should be so, and further, if \(\theta_{n}\) is the amplitude of \(1 + u_{n}\) (taken between \(-\pi\) and \(+\pi\) ), that the series of the \(\theta_{n}\) should be commutatively convergent. Since \(u_{n}\) tends to 0, \(\log(1 + u_{n})\) is defined starting from some particular value of n (III, p.~100) and one has

\[
\log (1 + u _ {n}) = \log | 1 + u _ {n} | + i \theta_ {n};
\]

thus, for the product with general factor \(1 + u_{n}\) to be commutatively convergent it is necessary and sufficient that the series with general term \(|\log(1 + u_{n})|\) be absolutely convergent (Gen.~Top., VII, p.~84, th. 1); now \(\log(1 + u_{n}) \sim u_{n}\) (I, p.~18, prop. 5), so one again obtains the condition that the series with general term \(u_{n}\) should be absolutely convergent (Gen.~Top., VIII, p.~116, th. 1).

The relation between infinite products and series of real numbers sometimes allows one to obtain an asymptotic expansion for the partial product \(p_{n} = \prod_{k=1}^{n}(1 + u_{k})\) ; it suffices to have an asymptotic expansion for the partial sum \(s_{n} = \sum_{k=1}^{n} \log(1 + u_{k})\) , then to expand \(p_{n} = \exp(s_{n})\) ; one is thus brought back to the two problems examined earlier (V, p.~238 and p.~226).

Example: Stirling's formula. Let us seek an asymptotic expansion for \(n!\) ; this leads us to expand \(s_n = \sum_{p=1}^{n} \log p\) , and then \(\exp(s_n)\) . The method of \(n^{\circ} 2\) gives successively

\[
s _ {n} = \sum_ {p = 1} ^ {n} \log p \sim \int_ {1} ^ {n} \log t d t = n \log n - n + 1
\]

then

\[
\log n - \int_ {n - 1} ^ {n} \log t d t = \log n - (n \log n - (n - 1) \log (n - 1) - 1) \sim \frac {1}{2 n}
\]

whence

\[
s _ {n} = n \log n - n + \frac {1}{2} \log n + o (\log n).
\]

Then

\[
\log n - \int_ {n - 1} ^ {n} \log t d t - \frac {1}{2} (\log n - \log (n - 1)) \sim - \frac {1}{1 2 n ^ {2}}
\]

whence

\[
s _ {n} = n \log n - n + \frac {1}{2} \log n + k + \frac {1}{1 2 n} + o _ {1} \left(\frac {1}{n}\right)
\]

(k constant)

and one finally deduces (V, p.~226)

\[
n! = e ^ {k} n ^ {n + 1 / 2} e ^ {- n} \left(1 + \frac {1}{1 2 n} + o _ {2} \left(\left(\frac {1}{n}\right)\right)\right).\tag{5}
\]

We shall show in VII, p.~322, that \(e^k = \sqrt{2\pi}\) . The formula (5) (with this value of \(k\) ) is called Stirling's formula. In the same way, for every real number \(a\) not an integer \(>0\) , one shows that

\[
(a + 1) (a + 2) \dots (a + n) \sim \mathrm{K} (a) n ^ {n + a + \frac {1}{2}} e ^ {- n}.\tag{6}
\]

We shall determine the function \(\mathrm{K}(a)\) too (VII, p.~18). From formulae (5) and (6) one derives in particular that

\[
\binom {a} {n} \sim (- 1) ^ {n}   \varphi (a)   n ^ {- a - 1}\tag{7}
\]

for every real number \(a\) not an integer \(> 0\) , where \(\varphi(a)\) is a function of \(a\) which will be specified in VII, p.~322.

